% LuaLaTeX, run "latexmk" in repository root
\documentclass[a4paper,11pt]{ltjsarticle}
\usepackage[margin=27mm]{geometry}
\usepackage{amsmath,amssymb,amsthm,mathtools,mathrsfs}
\usepackage[unicode,colorlinks=true,linkcolor=blue,urlcolor=blue]{hyperref}
\theoremstyle{plain}
\newtheorem{proposition}{命題}[section]
\newtheorem{lemma}[proposition]{補題}
\theoremstyle{definition}
\newtheorem{definition}[proposition]{定義}
\newtheorem{conjecture}[proposition]{予想}
\theoremstyle{remark}
\newtheorem{remark}[proposition]{注意}
\newcommand{\PP}{\mathbb P}
\newcommand{\RR}{\mathbb R}
\newcommand{\CC}{\mathbb C}
\newcommand{\T}{\mathbb T}
\DeclareMathOperator{\Ric}{Ric}
\title{一般型曲面のLagrangian torus fibration\\相対接着・分岐底空間・整数ホモロジー}
\author{}
\date{2026年11月20日 / GitHub移行用再構成草稿 v0.1}
\begin{document}
\maketitle

\begin{abstract}
滑らかな一般型超曲面$X_d\subset\PP^3$に対し、特定のKähler class上で
branched baseを持つLagrangian torus fibrationを作る問題を考える。
2026年10月27日付の30頁原PDF（第5稿）はrelative compactification、
Liouville collars、divisor-preserving Moser同定、baseの整数ホモロジーを扱う。
本ファイルはその論理構造を整理した再構成草稿であり、
原PDFの全証明を忠実に転記したものではない。
\end{abstract}
\section{一般型曲面と偏極}
\begin{proposition}
滑らかな次数$d$の超曲面$X_d\subset\PP^3$について
\[
 K_{X_d}\simeq\mathcal O_{X_d}(d-4)
\]
である。特に$d\ge5$では$K_{X_d}$はampleである。
\end{proposition}
\begin{proof}
adjunction formulaにより
$K_{X_d}=(K_{\PP^3}+X_d)|_{X_d}=
\mathcal O_{\PP^3}(d-4)|_{X_d}$。
$d-4>0$ならhyperplane classの正倍はampleである。
\end{proof}
\section{基底の分岐とfibers}
原PDFの構成は通常の閉曲面ではない自由辺なしの
compact branched two-complex $B_d$を基底とする。
regular fiberは$\T^2$、branch edgeでは
$S^1\times(S^1\vee S^1)$という特異fibre modelを用いる。
\begin{definition}
$X$のsymplectic form $\omega$に関し、
$f:X\to B$のregular locusでfiber $F$が
$\dim_\RR F=\frac12\dim_\RR X$かつ$\omega|_F=0$なら、
そのregular fibersをLagrangianと呼ぶ。
\end{definition}
\section{相対構成の検証単位}
原PDFで中心となる接着の入力は次の四段階である。
\begin{enumerate}
\item pants decomposition上で、boundary fibersまで管理したlocal modelを構成する。
\item 指定Morse foliationを保つfinite-volume Liouville collarをとる。
\item divisorを保つMoser identificationにより指定Kähler classと整合させる。
\item positive normal momentで切ったregion上のequivariant symplectic
transportを行い、fiberwise Lagrangian条件とbranching dataを照合する。
\end{enumerate}
これらは互いに独立の局所技法だが、境界での相対条件を同時に満たすことが
global theoremの鍵である。原PDF v5はその構成を与えたと主張するが、
本移行稿では各接着補題を完全な形で再掲しておらず、今回独立に監査していない。
\section{基底の位相と比較}
原稿は各marked pants treeに対して、base $B_d$から
単体の2-skeletonへのhomotopy equivalenceと整数係数のchain homotopy、
$H_2(B_d;\mathbb Z)$のbasisを与える。
具体的なchain matricesの忠実な再現には原PDFと旧source ZIPの再照合が必要である。
局所正則部分のintegral-affine structureはLagrangian fibrationの
action coordinatesから得られるが、branch/coreを含むstrict tropical atlas、
special Lagrangian性は証明したと主張しない。
\section{公表までの最優先事項}
第一にPDF v5の全補題とその証明を正規のLuaLaTeXソースとして移す。
第二にMoser deformationが保持するKähler classの許容領域と
boundary moment mapsを厳密に確認する。
第三にbranch edgeからcoreへ移るmonodromyとintegral homology
の全接着を独立に監査する。
これはKE degenerationから自動的に従うものではない。

\end{document}
